% The atom-tail reduction — written by Fable oracle R12
% For integration into proof.tex
% Cross-reference labels need alignment with main paper

\section{The atom-tail reduction}
\label{sec:atom-tail}

Throughout this section $X\ge 3$ is an integer parameter,
$\mathcal P_X=\{p\ \text{prime}:X<p\le2X\}$,
$\mathcal Z_p=\{r\in[0,p):p\mid b_r\}$, $Z(p)=|\mathcal Z_p|$,
\[
  K_X(m)=\#\{p\in\mathcal P_X:\ m\bmod p\in\mathcal Z_p\},
  \qquad
  \lambda_X=\sum_{p\in\mathcal P_X}\frac{Z(p)}p ,
\]
and $(K)_k=K(K-1)\cdots(K-k+1)$.  We write
$\omega_X(N)=\#\{p\in\mathcal P_X:p\mid N\}$ for the number of
window prime divisors of a nonzero integer $N$, so that
$\omega_X(N)\le\log|N|/\log X$.  We freely use the second
factorial moment $\sum_{m<X^2}(K_X(m))_2\le5X^2\lambda_X^2$
(Subsection~\ref{ssec:high-moment}) and the gap polynomials
$N_h$ of degree $3(h-1)$ with $N_h\not\equiv0\pmod p$ for
$h\le p$.

\subsection{The Cartier diagonal}

\begin{theorem}[Cartier diagonal]\label{thm:atomtail-cartier}
Let $F=(1-x-y)(1-z-w)-xyzw$.  For every prime $p$ and every
$0\le r\le p-1$,
\[
  b_r\;\equiv\;[x^ry^rz^rw^r]\,F^{p-1}\pmod p .
\]
\end{theorem}

\begin{proof}
By Straub's diagonal representation \cite{Straub2014},
$b_n=[x^ny^nz^nw^n]\,F^{-1}$ for all $n\ge0$.
Modulo $p$ one has $F^{-1}=F^{p-1}\cdot F^{-p}$ and, by the
Frobenius endomorphism, $F(x,y,z,w)^p\equiv F(x^p,y^p,z^p,w^p)$.
Hence
\[
 [x^r\!\cdots w^r]F^{-1}
 \equiv\sum_{b\ge0}[x^{r-pb}\!\cdots w^{r-pb}]F^{p-1}\cdot
 [x^b\!\cdots w^b]F^{-1}\pmod p .
\]
Since $\deg_xF^{p-1}=p-1$ and $0\le r\le p-1$, only $b=0$
contributes, and $[x^0\cdots w^0]F^{-1}=b_0=1$.
\end{proof}

\begin{corollary}[Palindrome for all fibers]
\label{cor:atomtail-palindrome-lambda}
For every prime $p \ge 5$, every $0 \le r \le p-1$,
and every $\lambda$:
$b_r(\lambda) \equiv b_{p-1-r}(\lambda) \pmod{p}$,
where $b_r(\lambda) = \sum_{k=0}^r \binom{r}{k}^2\binom{r+k}{k}^2 \lambda^k$.
\end{corollary}

\begin{proof}
The Cartier representation
$b_r(\lambda) \equiv [x^r \cdots w^r] F_\lambda^{p-1} \pmod{p}$
holds for $F_\lambda = (1{-}x{-}y)(1{-}z{-}w) - \lambda xyzw$
(Theorem~\ref{thm:atomtail-cartier} with general $\lambda$).
The torus involution $(x,y,z,w) \mapsto (x^{-1},y^{-1},z^{-1},w^{-1})$
sends $F_\lambda^{p-1}$ to $(xyzw)^{-2(p-1)} G^{p-1}$
where $G = (xyzw)^2 F_\lambda(x^{-1},\ldots)$
has degree~$2$ in each variable.
Extracting $[x^r \cdots w^r]$ from
$(xyzw)^{p-1} F_\lambda(x^{-1},\ldots)^{p-1}$ gives
$[x^{p-1-r} \cdots w^{p-1-r}] F_\lambda^{p-1}$,
and the two representations are equal modulo~$p$.
\end{proof}

\begin{remark}[Uniqueness of the Ap\'ery fiber]
\label{rem:atomtail-uniqueness}
The polynomial family $b_r(\lambda)$ satisfies a
\emph{fourth}-order recurrence in~$r$ for generic~$\lambda$
(Zeilberger), which degenerates to the second-order Ap\'ery
recurrence \emph{exactly} at $\lambda = 1$:
the defect $(r{+}1)^3 b_{r+1}(\lambda) - P(r) b_r(\lambda)
+ r^3 b_{r-1}(\lambda)$
factors as $C_r \cdot (\lambda - 1) \cdot Q_r(\lambda)$
with $\gcd(Q_1, \ldots, Q_7) = 1$.
Thus $\lambda = 1$ is the unique rational specialization
admitting the order-two certificate algebra
(gap polynomials, codegree bounds, per-prime separation)
on which every zero-set theorem in this paper rests.
\end{remark}

\subsection{Fourier non-concentration}

Recall $F_p(a)=\sum_{r\in\mathcal Z_p}e_p(ar)$ and set
$A_p(h)=\#\{r:\ r,r+h\in\mathcal Z_p\}$ (integer gap,
$0\le r<r+h\le p-1$), so that $A_p(h)\le3(h-1)$: indeed
$r,r+h\in\mathcal Z_p$ forces $N_h(r)\equiv0\pmod p$, and
$N_h\not\equiv0\pmod p$ has at most $\deg N_h=3(h-1)$ roots.

\begin{lemma}[Short-window energy]\label{lem:atomtail-energy}
There is an absolute constant $C$ such that for every prime $p$
and every $1\le K_0\le p/2$,
\[
 \sum_{1\le|k|\le K_0}|F_p(k)|^2
 \;\le\;4K_0\,Z(p)\;+\;C\,\frac{p^2\log(2K_0)}{K_0},
\]
while for $K_0>p/2$ the full Parseval bound
$\sum_{k\bmod p}|F_p(k)|^2=pZ(p)-Z(p)^2$ applies.
\end{lemma}

\begin{proof}
Let $w(k)=(1-|k|/(2K_0))_+$ on the symmetric residue window;
$w\ge\frac12$ on $|k|\le K_0$, so the left side is at most
$2\sum_{k\bmod p}w(k)|F_p(k)|^2
=2\sum_{r,r'\in\mathcal Z_p}\widehat W(r-r')$, where
\[
 \widehat W(h)=\frac1{2K_0}\Bigl|\sum_{0\le j<2K_0}e_p(jh)\Bigr|^2
 \;\le\;\min\Bigl(2K_0,\ \frac{p^2}{8K_0\bar h^2}\Bigr),
 \qquad \bar h=\min(|h|,p-|h|).
\]
The diagonal contributes $2K_0Z(p)$.  Off-diagonal pairs at
integer gap $g$ number $A_p(g)\le3(g-1)$ for $g\le p/2$, and at
most $p-g$ for $g>p/2$ (the smaller element lies in
$[0,p-1-g]$); in either case the count at circular distance
$\bar h$ is at most $3\bar h$.  Splitting at
$g^*=p/(4K_0)$,
\[
 \sum_{\bar h\ge1}3\bar h\,
 \min\Bigl(2K_0,\frac{p^2}{8K_0\bar h^2}\Bigr)
 \le 6K_0\!\!\sum_{\bar h\le g^*}\!\!\bar h
 +\frac{3p^2}{8K_0}\!\!\sum_{g^*<\bar h\le p/2}\!\!\frac1{\bar h}
 \;\ll\;\frac{p^2\log(2K_0)}{K_0}. \qedhere
\]
\end{proof}

\subsection{The codegree dichotomy and the twin-atom bound}

For $0\le m<m'<X^2$ put
$W(m,m')=\#\{p\in\mathcal P_X:\ m\bmod p\in\mathcal Z_p,\
m'\bmod p\in\mathcal Z_p\}$ and $h=m'-m$.

\begin{lemma}[Codegree dichotomy]\label{lem:atomtail-dichotomy}
Let $1\le h<X$ and let $p\in\mathcal P_X$ be counted by
$W(m,m')$, with $r=m\bmod p$.  Exactly one of the following
holds.
\begin{enumerate}
\item[\textup{(A)}] $r+h\le p-1$.  Then $r,r+h\in\mathcal Z_p$,
hence $N_h(r)\equiv0\pmod p$ and, since $N_h\in\mathbb Z[x]$,
\[
  p\mid N_h(m).
\]
\item[\textup{(B)}] $r\ge p-h$.  Then $j:=m'\bmod p=r+h-p$
satisfies $0\le j<h$, $j\in\mathcal Z_p$ means $p\mid b_j$, and
$p\mid m'-j$; hence
\[
  p\mid\gcd\bigl(b_j,\ m'-j\bigr)
  \quad\text{for some }0\le j<h.
\]
\end{enumerate}
\end{lemma}

\begin{proof}
Since $h<X<p$, the residues satisfy $m'\bmod p\equiv r+h$ with
$r+h<2p$, giving the two cases.  In case (A) both indices lie in
$[0,p-1]$ and the gap certificate applies; the congruence
$N_h(m)\equiv N_h(r)\pmod p$ is the reduction of a polynomial
with integer coefficients.  Case (B) is immediate from the
definitions.
\end{proof}

\begin{lemma}[Positivity of gap values]\label{lem:atomtail-positive}
$N_h(m)\ge1$ for all integers $m\ge0$, $h\ge1$.
\end{lemma}

\begin{proof}
$N_h(m)=\det D_h(m)$ with $D_h(m)$ symmetric tridiagonal,
diagonal $P(m+i)$ and off-diagonal $(m+i+1)^3$.
The identity
\[
  P(n)-n^3-(n+1)^3=4(2n+1)^3>0\qquad(n\ge0)
\]
shows $D_h(m)$ is strictly diagonally dominant with positive
diagonal, hence positive definite.
\end{proof}

\begin{theorem}[Twin-atom bound; repaired codegree lemma]
\label{thm:atomtail-twin}
For $0\le m<m'<X^2$ with $h=m'-m<X$,
\[
 W(m,m')\;\le\;\omega_X\bigl(N_h(m)\bigr)
 +\sum_{j=0}^{h-1}\omega_X\bigl(\gcd(b_j,m'-j)\bigr)
 \;\le\;7h+\min\Bigl(2h,\ \frac{2h^2}{\log X}\Bigr),
\]
and $W(m,m+1)=0$ identically.
\end{theorem}

\begin{proof}
The first inequality is Lemma~\ref{lem:atomtail-dichotomy}
combined with Lemma~\ref{lem:atomtail-positive}.  For the
numerical form: the entries of $D_h(m)$ are at most $C X^6$, so
$\log N_h(m)\le 6(h-1)\log X+O(h\log h)$ and
$\omega_X(N_h(m))\le7h$ for large $X$; each $m'-j<X^2$ has at
most two window prime divisors, and
$\omega_X(b_j)\le\log b_j/\log X\le 3.6j/\log X$ gives the
alternative bound.  For $h=1$: case (A) requires
$N_1\equiv0$ with $N_1=1$, and case (B) requires $p\mid b_0=1$;
both are impossible.
\end{proof}

\subsection{Short-range pair statistics}

\begin{proposition}[Short-range decorrelation]
\label{prop:atomtail-shortrange}
For $1\le h<X$,
\begin{align}
 \sum_{m<X^2-h}W(m,m+h)
 &\;\le\;C\,\frac{h\,X^2}{\log X},
 \label{eq:atomtail-t1-first}\\
 \sum_{m<X^2-h}W(m,m+h)\bigl(W(m,m+h)-1\bigr)
 &\;\le\;C\,\frac{h^2X^2}{(\log X)^2}.
 \label{eq:atomtail-t1-energy}
\end{align}
\end{proposition}

\begin{proof}
For \eqref{eq:atomtail-t1-first}: type-(A) incidences are at most
$\sum_{p}\#\{m<X^2:p\mid N_h(m)\}
\le\sum_p 3(h-1)(X^2/p+1)\ll hX^2/\log X$ by Mertens; type-(B)
incidences are at most
$\sum_{j<h}\sum_{p\mid b_j}(X^2/p+1)
\le\sum_{j<h}\omega_X(b_j)\,(X+1)\ll h^2X/\log X$,
which is smaller.  For \eqref{eq:atomtail-t1-energy}: an ordered
pair $p\ne q$ counted together is, in the dominant (A,A)
configuration, a pair with $pq\mid N_h(m)$, and
\[
 \sum_{p\ne q}\#\{m<X^2:pq\mid N_h(m)\}
 \le\sum_{p\ne q}9h^2\Bigl(\frac{X^2}{pq}+1\Bigr)
 \ll\frac{h^2X^2}{(\log X)^2};
\]
mixed configurations are bounded by the type-(B) count times the
pointwise bound of Theorem~\ref{thm:atomtail-twin}.
\end{proof}

\subsection{The root-strip theorem and resultant nonvanishing}

\begin{theorem}[Root strip]\label{thm:atomtail-rootstrip}
For every $h \ge 2$, all roots of $N_h$ lie in the open strip
$\{z \in \mathbb{C} : -h < \operatorname{Re} z < -1\}$.
\end{theorem}

\begin{proof}
\emph{Step~1.}  For $\operatorname{Re} y \ge 0$, writing
$u = |y + \tfrac12| \ge \tfrac12$:
\[
  |P(y)| = 2u\,|17(y{+}\tfrac12)^2 + \tfrac34|
  \ge 2u(17u^2 - \tfrac34) = 34u^3 - \tfrac32 u,
\]
while $|y|^3 + |y{+}1|^3 \le 2(u + \tfrac12)^3
= 2u^3 + 3u^2 + \tfrac32 u + \tfrac14$.
The difference $32u^3 - 3u^2 - 3u - \tfrac14$ equals
$\tfrac32$ at $u = \tfrac12$ and has positive derivative
$96u^2 - 6u - 3 > 0$ for $u \ge \tfrac12$.
Hence $|P(y)| > |y|^3 + |y{+}1|^3$.

\emph{Step~2.}  $N_h(z) = \det D_h(z)$ with
$D_h$ tridiagonal: diagonal $P(z{+}i)$,
off-diagonal $(z{+}i{+}1)^3$.
For $\operatorname{Re} z \ge -1$, every $y_i = z + i$
($1 \le i \le h{-}1$) has $\operatorname{Re} y_i \ge 0$,
so Step~1 gives strict diagonal dominance
(endpoint rows have only one off-diagonal entry, strengthening
the inequality).
By Levy--Desplanques, $\det D_h(z) \ne 0$.

\emph{Step~3.}  The reflection identity
$N_h(-z{-}h{-}1) = (-1)^{h-1} N_h(z)$
(Remark~\ref{rem:dodgson}, valid for all $z \in \mathbb{C}$
by the polynomial identity of the tridiagonal determinant)
maps any root $\alpha$ to the root $-\alpha{-}h{-}1$
with $\operatorname{Re}(-\alpha{-}h{-}1)
= -\operatorname{Re}\alpha - h - 1$.
If $\operatorname{Re}\alpha \le -h$, then
$\operatorname{Re}(-\alpha{-}h{-}1) \ge -1$,
contradicting Step~2.
\end{proof}

\begin{corollary}[Resultant nonvanishing --- the First Lemma]
\label{cor:atomtail-firstlemma}
For all $d, r \ge 1$,
$S_{d,r} = \operatorname{Res}_x(N_d(x),\, N_r(x{+}d)) \ne 0$.
For $d = 1$ or $r = 1$, this is immediate since $N_1 = 1$.
For $d, r \ge 2$: by Theorem~\ref{thm:atomtail-rootstrip},
roots of $N_d$ have $\operatorname{Re} \in (-d, -1)$
while roots of $N_r(x{+}d)$
have $\operatorname{Re} \in (-d{-}r,\, -d{-}1)$.
The two strips are disjoint (separated by $[-d{-}1, -d]$),
so the polynomials share no root and their resultant is nonzero.
\end{corollary}

\subsection{Per-prime separation}

Recall $S_{h_1,h_2}=\operatorname{Res}_x\bigl(N_{h_1}(x),\,
N_{h_2}(x+h_1)\bigr)\in\mathbb Z$, with
$\log|S_{h_1,h_2}|\ll(h_1+h_2)^2\log(h_1+h_2)$.
By Corollary~\ref{cor:atomtail-firstlemma},
$S_{h_1,h_2} \ne 0$ unconditionally.

\begin{theorem}[Per-prime separation]\label{thm:atomtail-sep}
Let $h_1,h_2\ge1$.  If some
$p\in\mathcal P_X$ satisfies
$m,\,m+h_1,\,m+h_1+h_2\in\Omega_p$ (all three hit by $p$), then
\[
 p\mid S_{h_1,h_2}
 \qquad\text{or}\qquad
 p\mid b_j\ \text{for some }0\le j<h_1+h_2 .
\]
(The hypothesis $S_{h_1,h_2}\ne0$ from earlier versions of this
result is now unconditional by
Corollary~\ref{cor:atomtail-firstlemma}.)
Consequently there is $c>0$ such that for
$h_1+h_2\le c\,(\log X/\log\log X)^{1/2}$ and $X$ large, no
prime of $\mathcal P_X$ hits all three points:  each hit set
$\Omega_p$ contains at most two points in any window of length
$c\,(\log X/\log\log X)^{1/2}$.
\end{theorem}

\begin{proof}
If neither consecutive pair wraps
(Lemma~\ref{lem:atomtail-dichotomy}(A) twice), then with
$r=m\bmod p$ both $N_{h_1}(r)\equiv0$ and
$N_{h_2}(r+h_1)\equiv0\pmod p$.  Since $p>X>h_1+h_2$, the
leading coefficients of $N_{h_1}$ and $N_{h_2}(\cdot+h_1)$ are
units modulo $p$, so a common root forces
$p\mid\operatorname{Res}=S_{h_1,h_2}$.  If some pair wraps, case
(B) gives $p\mid b_j$ with $j<h_1+h_2$.  For the quantitative
form: $\omega_X(S_{h_1,h_2})\le
C(h_1+h_2)^2\log(h_1+h_2)/\log X<1$ in the stated range, so no
window prime divides $S_{h_1,h_2}$; and $b_j\le
e^{3.6j}<X<p$ in the stated range, so $p\mid b_j$ is impossible.
\end{proof}

\begin{remark}
The nonvanishing $S_{h_1,h_2} \ne 0$ is now unconditional
(Corollary~\ref{cor:atomtail-firstlemma}).
For consecutive indices the stronger product formula
$|\operatorname{Res}(N_h,N_{h+1})|=\prod_j(j!^3b_j)^6$
gives the exact value.
The root-strip theorem
(Theorem~\ref{thm:atomtail-rootstrip})
provides the sharpest known localization:
every root of $N_h$ lies on the ``critical strip''
$-h < \operatorname{Re} z < -1$,
centered on the reflection line
$\operatorname{Re} z = -(h{+}1)/2$.
\end{remark}

\subsection{The moment collapse and the atom-tail equivalence}

\begin{proposition}[Moment collapse]\label{prop:atomtail-collapse}
For every $k\ge3$,
\[
 \sum_{m<X^2}(K_X(m))_k
 \;\le\;\Bigl(\max_{m<X^2}K_X(m)\Bigr)^{k-2}
 \sum_{m<X^2}(K_X(m))_2
 \;\le\;5X^2\lambda_X^2\,\Bigl(\max_{m<X^2}K_X(m)\Bigr)^{k-2}.
\]
\end{proposition}

\begin{proof}
Pointwise, $(K)_k=(K)_2\prod_{j=2}^{k-1}(K-j)\le(K)_2K^{k-2}$;
sum over $m$ and apply the second factorial moment.
\end{proof}

\begin{hypothesis}[Atom bound $(\mathrm{AT}'')$]
\label{hyp:atomtail-at}
$\displaystyle\max_{m<X^2}K_X(m)\;\le\;\lambda_X\,X^{o(1)}.$
\end{hypothesis}

\begin{corollary}\label{cor:atomtail-at}
$(\mathrm{AT}'')$ implies $(\mathrm{HM})_k$ for every fixed
$k\ge2$ simultaneously; in particular, by
Theorem~\ref{thm:hm-pointwise} with $k=7$,
$(\mathrm{AT}'')$ implies
$\log G_n\ll n^{20/21+o(1)}$ for every $n$, hence
$G_n=e^{o(n)}$.
Combined with the unconditional average bound
$\sum_{p \le x} Z(p) \ll x^{8/5}/\log x$
(from Corollary~\ref{cor:atomtail-firstlemma} and
the corank-valuation argument of
Subsection~\ref{ssec:value-fibers}),
one obtains $\lambda_X \ll X^{3/5}/\log X$, so
$(\mathrm{AT}'')$ with $k = 6$ gives
$\log G_n \ll n^{14/15 + o(1)}$.
\end{corollary}

\begin{proof}
Proposition~\ref{prop:atomtail-collapse} gives
$\sum(K)_k\le5X^2\lambda_X^2(\lambda_XX^{o(1)})^{k-2}
=X^{2+o(1)}\lambda_X^k$.
\end{proof}

\begin{theorem}[Concentration of the problem in the atom tail]
\label{thm:atomtail-conc}
Let $n_T(X)=\#\{m<X^2:K_X(m)\ge T\}$.  Then:
\begin{enumerate}
\item Unconditionally,
$n_T(X)\le5X^2\lambda_X^2/(T)_2$ for all $T\ge2$, and
$\max_mK_X(m)\le|\mathcal P_X|\ll X/\log X$.
\item $(\mathrm{HM})_3$ holds if and only if
\[
 n_T(X)\;\ll\;\frac{X^{2+o(1)}\lambda_X^3}{T^3}
 \qquad\text{uniformly for }T\ge3 .
\]
\item If $n_T(X)\ll X^{2+o(1)}\lambda_X^k/T^k$ uniformly for
$T\ge k$, then $(\mathrm{HM})_k$ holds; with $k=7$ this implies
$G_n=e^{o(n)}$ for every $n$.
\end{enumerate}
\end{theorem}

\begin{proof}
(1) is Chebyshev applied to the second factorial moment, and the
trivial ceiling.  (2), necessity: $n_T\,(T)_3\le\sum_m(K)_3$.
Sufficiency: writing the third factorial moment as a sum over
levels,
\[
 \sum_m(K_X(m))_3
 =\sum_{T\ge3}n_T\bigl[(T)_3-(T-1)_3\bigr]
 \le\sum_{T\ge3}3T^2\,n_T
 \ll X^{2+o(1)}\lambda_X^3\sum_{3\le T\ll X}\frac1T,
\]
and the harmonic sum is $O(\log X)=X^{o(1)}$.  (3) is identical
with $[(T)_k-(T-1)_k]\le kT^{k-1}$, followed by
Theorem~\ref{thm:hm-pointwise}.
\end{proof}

\begin{remark}[Status of the frontier]
\label{rem:atomtail-frontier}
The gap between item (1) and item (2) is a single factor
$T/\lambda_X$ per level: Chebyshev gives
$X^2\lambda_X^2/T^2$, the conjecture requires
$X^2\lambda_X^3/T^3$.  The Poisson model predicts
$n_T\approx X^2\lambda_X^T/T!$ and
$\max K_X\asymp\log X/\log\log X$, in agreement with
computation ($\max K=4$ at $X=1024$; nearest-neighbour
statistics of the level sets Poisson to three digits).  Every
ensemble method examined in this program---large-sieve
dispersion, Linnik dispersion with completed CRT sums, variance
layer-cake, codegree-weighted pair dispersion---either requires
the atom bound at its diagonal or is subsumed by
Proposition~\ref{prop:atomtail-collapse}.  The full conjecture
is thereby reduced to a single pointwise statement about the
divisor structure of the Ap\'ery sequence: no integer
$m<X^2$ may lie in the zero fibre of
$X^{\varepsilon}$ distinct window primes.
\end{remark}

\subsection{The corank ladder and its ceiling}

The average exponent $c$ in $\sum_{p \le x} Z(p) \ll x^{1+c}/\log x$
determines the moment order needed for the full conjecture via the
trade-off $c + 2/k < 1$.

\begin{center}
\begin{tabular}{ccc}
\toprule
$c$ & $k$ needed & status \\
\midrule
$2/3$ & $7$ & proved (gap polynomial) \\
$3/5$ & $6$ & proved (corank + First Lemma) \\
$1/2$ & 5 & conditional on $\mathrm{AC}^{\mathrm{tr}}$,
$\mathrm{FDAC}^{\mathrm{tr}}$, \textup{(SWR)}, or \textup{(QPRS)}
(\S\ref{ssec:quadcorank}) \\
$o(1)$ & $3$ & requires parity-breaking \\
\bottomrule
\end{tabular}
\end{center}

The pair corank argument (Gershgorin root-strip,
Corollary~\ref{cor:atomtail-firstlemma}) gives $c = 3/5$.
The na\"ive adjacent-pair quadruple certificate
$\gcd(S_{h_1,h_2},S_{h_2,h_3})$ is not transverse: when
$h_1=h_3$, the symmetry
$|S_{h_1,h_2}|=|S_{h_2,h_1}|$ makes it a self-gcd.  In
particular, the previously reported numerical decay for this mass
had inadvertently omitted its diagonal.  The skipped triple formed by
the first, second, and fourth zeros instead pairs
$S_{h_1,h_2}$ with $S_{h_1,h_2+h_3}$.  This removes the literal
self-gcd, but a second degeneracy remains when $h_3=h_1$: if $h_2$
is even, the central factor of $N_{h_2}$ propagates into both
resultants (Proposition~\ref{prop:meso-center-recurrence}).  The
transverse formulation below deflates that linear factor.  It then
replaces the scalar gcd of two resultants by the coefficient content of
one pencil resultant.  This aligned content records that the \emph{same}
starting point is a root of all three gap polynomials, rather than merely
that the same prime divides two resultants for possibly unrelated roots.
The still more singular arithmetic-progression case
$h_1=h_2=h_3$ has only one gap parameter and is handled directly by a
single resultant.
The separate obstruction at the triple level is that
$U_{h_1,h_2}
= \operatorname{Res}(N_{h_1}, N_{h_1+h_2})$ factors as
$U = \pm S_{h_1,h_2} \cdot \operatorname{Res}(N_{h_1}, N_{h_1+1})
\cdot (\text{leading-coefficient})$,
so certificates within one triple are algebraically dependent.

\subsection{The quadruple corank argument: average $Z(p)\ll p^{1/2}$}
\label{ssec:quadcorank}

Let $\varepsilon=17+12\sqrt2$ and
$V_n=(\varepsilon^n-\varepsilon^{-n})/(\varepsilon-\varepsilon^{-1})$
the Lucas sequence with $V_{n+1}=34V_n-V_{n-1}$, $V_0=0$, $V_1=1$,
so that $V_{n}=U_{n-1}(17)$ is the leading coefficient of $N_{n}$.

\begin{definition}[Structural content and reduced resultants]
For $H\ge2$ put
$\mathfrak U_H=\prod_{j\le H} j!\cdot b_j\cdot V_{j+1}$,
and define the saturated reduced resultant
$S^*_{h_1,h_2}$ as $S_{h_1,h_2}$ with all content shared with
$\mathfrak U_H$ removed:
$S^*=S/\gcd(S,\mathfrak U_H^\infty)$.
For $b\ge2$ put
\[
 N_b^\circ(x)=
 \begin{cases}
  N_b(x)/(2x+b+1),&b\ \text{even},\\
  N_b(x),&b\ \text{odd},
 \end{cases}
 \qquad
 D_{a,b}=\operatorname{Res}_x
 \bigl(N_a(x),N_b^\circ(x+a)\bigr),
\]
and define $D^*_{a,b}$ by the same saturation.
\end{definition}

\begin{lemma}[Skipped-triple quadruple certificate]\label{lem:qc-cert}
If $p$ has four zeros
$r<r{+}h_1<r{+}h_1{+}h_2<r{+}h_1{+}h_2{+}h_3\le p{-}1$
in $\mathcal Z_p$, then
$p\mid\gcd(S_{h_1,h_2},\,S_{h_1,h_2+h_3})$.
\end{lemma}

\begin{proof}
No two of the displayed zeros are consecutive
(Lemma~\ref{lem:no-consec}), so $h_i\ge2$.  Applying the gap-polynomial
identity to the first, second, and third zeros, and then to the first,
second, and fourth zeros, gives
\[
 N_{h_1}(r)=N_{h_2}(r+h_1)
 =N_{h_2+h_3}(r+h_1)=0\pmod p.
\]
There is no denominator problem because every index involved lies in
$[0,p-1]$.  Thus $r$ is a common root of
$N_{h_1}(x)$ and $N_{h_2}(x+h_1)$, and also of
$N_{h_1}(x)$ and $N_{h_2+h_3}(x+h_1)$.  Each common root makes the
corresponding Sylvester determinant vanish modulo~$p$.  The two
determinants are the nonzero integers $S_{h_1,h_2}$ and
$S_{h_1,h_2+h_3}$ by Corollary~\ref{cor:atomtail-firstlemma}, proving the
claim.
\end{proof}

\begin{lemma}[Central deflation]\label{lem:qc-deflation}
The polynomial $N_b^\circ$ belongs to $\mathbb Z[x]$, and
$D_{a,b}\ne0$ for all $a,b\ge2$.  Consider any partition of
$[0,p-1]$, for an odd prime~$p$, into intervals and select
consecutive quadruples of zeros
inside those intervals.  Among all selected quadruples with
palindromic gap vector $(a,b,a)$, at most one has
\begin{equation}\label{eq:qc-central-start}
 b\ \text{even},\qquad 2(r+a)+b+1\equiv0\pmod p.
\end{equation}
Every other such quadruple has $r$ as a common root modulo~$p$ of
$N_a(x)$ and $N_b^\circ(x+a)$, as well as of $N_a(x)$ and
$N_{a+b}(x+a)$.
\end{lemma}

\begin{proof}
When $b$ is even, the reflection identity
$N_b(-x-b-1)=(-1)^{b-1}N_b(x)$ makes
$-(b+1)/2$ a root of~$N_b$.  The primitive polynomial
$2x+b+1$ therefore divides $N_b$ in $\mathbb Z[x]$ by Gauss's
lemma.  Every root of $N_b^\circ(x+a)$ is a root of
$N_b(x+a)$, so the root-strip theorem gives $D_{a,b}\ne0$.

For a palindromic quadruple, the gap identities give
\[
 N_a(r)=N_b(r+a)=N_{a+b}(r+a)=0\pmod p.
\]
Unless~\eqref{eq:qc-central-start} holds, the factor
$2(r+a)+b+1$ is a unit, so $N_b^\circ(r+a)=0$ as well.
If~\eqref{eq:qc-central-start} does hold, then its endpoints satisfy
\[
 2r+2a+b\equiv-1\pmod p.
\]
Both endpoints lie in $[0,p-1]$, hence
$r+(r+2a+b)=p-1$.  By
Corollary~\ref{cor:atomtail-palindrome-lambda}, the Ap\'ery zero set is
invariant under $z\mapsto p-1-z$.  If its elements are
$u_1<\cdots<u_{Z(p)}$, a consecutive quadruple
$u_i,\ldots,u_{i+3}$ with $u_i+u_{i+3}=p-1$ must satisfy
$i=Z(p)-i-2$, which determines~$i$ uniquely.  A quadruple
consecutive inside one partition interval is also consecutive in the
full zero set, proving the asserted global uniqueness.
\end{proof}

\begin{definition}[Aligned pencil content]\label{def:qc-aligned}
For $a,b,c\ge2$ that are not all equal, define three
polynomials in $\mathbb Z[x]$ by
\[
 (F_{a,b,c},G_{a,b,c},J_{a,b,c})=
 \begin{cases}
 \bigl(N_a(x),N_b(x+a),N_{b+c}(x+a)\bigr),&c\ne a,\\
 \bigl(N_a(x),N_b^\circ(x+a),N_{a+b}(x+a)\bigr),
     &c=a,\ a\ne b.
 \end{cases}
\]
Put
\[
 \mathcal R_{a,b,c}(T)
  =\operatorname{Res}_x\bigl(F_{a,b,c},G_{a,b,c}+T J_{a,b,c}\bigr)
  \in\mathbb Z[T],
 \qquad
 \mathcal C_{a,b,c}=\operatorname{cont}_T\mathcal R_{a,b,c},
\]
where $\operatorname{cont}_T$ denotes the positive gcd of all integer
coefficients.  Relative to an ambient height~$H$, let
$\mathcal C^*_{a,b,c}$ be obtained from $\mathcal C_{a,b,c}$ by
removing its full prime-power part supported on~$\mathfrak U_H$.
\end{definition}

\begin{lemma}[Aligned-content valuation]\label{lem:qc-aligned}
Let $F,G,J\in\mathbb Z[x]$ have positive generic degrees, let
\[
 \mathcal R(T)=\operatorname{Res}_x(F,G+TJ)\ne0,
 \qquad \mathcal C=\operatorname{cont}_T\mathcal R,
\]
and suppose that $\overline F\ne0$ in $\mathbb F_p[x]$.  If
$\overline F,\overline G,\overline J$ have $t$ distinct common roots
in~$\mathbb F_p$, then
\begin{equation}\label{eq:qc-aligned-valuation}
 t\le v_p(\mathcal C).
\end{equation}
No simplicity assumption on those roots is required.

Moreover, every pencil in Definition~\ref{def:qc-aligned} is nonzero,
and after saturation one has
\begin{equation}\label{eq:qc-aligned-divisibility}
 \mathcal C^*_{a,b,c}\mid
 \begin{cases}
  \gcd(S^*_{a,b},S^*_{a,b+c}),&c\ne a,\\
  \gcd(D^*_{a,b},S^*_{a,a+b}),&c=a,\ a\ne b.
 \end{cases}
\end{equation}
\end{lemma}

\begin{proof}
Equip $K=\mathbb Q_p(T)$ with the Gauss valuation
\[
 v_{\mathrm G}\!\left(\sum_jq_jT^j\right)=\min_jv_p(q_j)
\]
and extend it to rational functions.  Its valuation ring is a discrete
valuation ring with uniformizer~$p$ and residue field
$k=\mathbb F_p(T)$.  Write $m=\deg F$ and
$n=\max(\deg G,\deg J)$.  Over this valuation ring the generic
Sylvester map is
\[
 \Phi:\mathcal O[x]_{<n}\oplus\mathcal O[x]_{<m}
 \longrightarrow\mathcal O[x]_{<m+n},
 \qquad (A,B)\longmapsto AF+B(G+TJ),
\]
and its determinant is $\pm\mathcal R(T)$.

Let $\alpha_1,\ldots,\alpha_t$ be the common roots in~$\mathbb F_p$.
Evaluation at each~$\alpha_i$ defines a $k$-linear functional on the
target of $\overline\Phi$, and all these functionals annihilate its
image.  They are linearly independent: since $\overline F\ne0$, one
has $t\le\deg\overline F\le m$, and Lagrange interpolation separates
the $t$ evaluation maps on $k[x]_{<m+n}$.  Thus
$\overline\Phi$ has corank at least~$t$.  Smith normal form over the
Gauss valuation ring gives
\[
 v_{\mathrm G}(\det\Phi)\ge t.
\]
If $\mathcal R(T)=\sum_jr_jT^j$, then
\[
 v_{\mathrm G}(\mathcal R)=\min_jv_p(r_j)
 =v_p\!\left(\gcd_jr_j\right)=v_p(\mathcal C),
\]
which proves~\eqref{eq:qc-aligned-valuation}.  Notice that the
Sylvester matrix has its generic size before reduction; a drop of a
leading coefficient modulo~$p$ cannot invalidate this corank argument.

For the pencils in Definition~\ref{def:qc-aligned}, put
$d_F=\deg F_{a,b,c}$, $d_G=\deg G_{a,b,c}$, and
$d_J=\deg J_{a,b,c}$.  In both cases $d_J>d_G$, and specialization of
the generic-degree resultant gives
\[
 \mathcal R_{a,b,c}(0)
 =\operatorname{lc}(F_{a,b,c})^{d_J-d_G}
   \operatorname{Res}_x(F_{a,b,c},G_{a,b,c}),
\]
whereas the coefficient of $T^{d_F}$ is
$\operatorname{Res}_x(F_{a,b,c},J_{a,b,c})$.  The former resultant is
$S_{a,b}$ off the palindromic slice and $D_{a,b}$ on that slice, up to
sign; the latter is $S_{a,b+c}$, again up to sign.  They are nonzero by
Corollary~\ref{cor:atomtail-firstlemma} and
Lemma~\ref{lem:qc-deflation}.  Hence the pencil resultant is nonzero.
Its content divides both displayed endpoint coefficients.  Finally,
$\operatorname{lc}(F_{a,b,c})=V_a$, whose prime support is included in
$\mathfrak U_H$.  Removing all $\mathfrak U_H$-supported prime powers
therefore gives~\eqref{eq:qc-aligned-divisibility}.
\end{proof}

\begin{proposition}[Normalization length of pencil content]
\label{prop:qc-normalization-length}
Let $R$ be a complete discrete valuation ring of characteristic zero
with perfect residue field, fraction field~$K$, and valuation~$v_R$.
Let $F,G,J\in R[x]$,
assume that the leading coefficient of~$F$ is a unit, and suppose
$\operatorname{Res}_x(F,G+TJ)\ne0$.  Let $\widehat F$ be the monic
polynomial obtained from~$F$ by dividing by its leading coefficient, and
factor it over~$K$ as
\[
 \widehat F=\prod_i f_i^{m_i},
\]
where the $f_i$ are distinct monic irreducibles.  Put
$L_i=K[x]/(f_i)$, let $\mathcal O_i$ be the valuation ring of~$L_i$,
and write $\alpha_i$ for the image of~$x$.  Then
\begin{equation}\label{eq:qc-normalization-length}
 v_R\!\left(\operatorname{cont}_T
       \operatorname{Res}_x(F,G+TJ)\right)
 =\sum_i m_i\,
   \ell_R\!\left(
    \mathcal O_i/(G(\alpha_i),J(\alpha_i))
   \right).
\end{equation}
\end{proposition}

\begin{proof}
The unit leading coefficient makes no contribution to the valuation.
Multiplicativity of the resultant gives, up to a unit of~$R$,
\[
 \operatorname{Res}_x(F,G+TJ)
 =\prod_i
   \operatorname N_{L_i/K}
       \bigl(G(\alpha_i)+T J(\alpha_i)\bigr)^{m_i}.
\]
Let $e_i$ and $f_i^{\mathrm{res}}$ denote the ramification index and
residue degree of $L_i/K$, and normalize the integral valuation
$v_i$ on~$L_i$ by $v_i|_K=e_i v_R$.  The Gauss valuation is
multiplicative.  After passing to an algebraic closure of~$K$, every
conjugate of the linear polynomial
$G(\alpha_i)+TJ(\alpha_i)$ has Gauss valuation
\[
 \frac1{e_i}\min\{v_i(G(\alpha_i)),v_i(J(\alpha_i))\}.
\]
There are $e_i f_i^{\mathrm{res}}$ conjugates, so
\[
 v_R\!\left(
  \operatorname N_{L_i/K}(G(\alpha_i)+TJ(\alpha_i))
 \right)
 =f_i^{\mathrm{res}}
  \min\{v_i(G(\alpha_i)),v_i(J(\alpha_i))\}.
\]
The right-hand side is exactly the $R$-length of
$\mathcal O_i/(G(\alpha_i),J(\alpha_i))$.  Summing over the branches
with multiplicities~$m_i$ proves~\eqref{eq:qc-normalization-length},
because the Gauss valuation of a polynomial over~$K$ is the minimum
valuation of its coefficients.
\end{proof}

\begin{remark}
Formula~\eqref{eq:qc-normalization-length} is a length on the
normalization of the reduced $K$-algebra defined by~$F$, with the
scheme-theoretic multiplicities of~$F$ restored.  It is not, in
general, the length of $R[x]/(F,G,J)$: the order $R[x]/(F)$ need not
be integrally closed, and normalization can change the local length.
The two descriptions agree, for example, when $R[x]/(F)$ is finite
\'{e}tale over~$R$.  A concrete discrepancy is obtained, for odd
residue characteristic, from
\[
 F=x^2-\varpi^2,\qquad G=x-\varpi,\qquad J=x+\varpi.
\]
Here
$\operatorname{Res}_x(F,G+TJ)=-4\varpi^2T$, so the content valuation
and the sum of the two normalized branch lengths are both~$2$,
whereas $R[x]/(F,G,J)\cong R/(2\varpi)$ has length~$1$.
\end{remark}

\begin{lemma}[Rank of apparition]\label{lem:qc-rank}
For $p\nmid6$, put $\chi(p)=(2/p)$.  Then
$p\mid V_n\iff\rho(p)\mid n$ for an integer
$\rho(p)\mid p-\chi(p)$.  Hence
$\#\{p\in\mathcal P_X:\rho(p)\le R\}\ll R^2/\log X$.
\end{lemma}

\begin{proof}
Work in $\mathbb F_p(\sqrt2)$ and write
$\bar\varepsilon=17+12\sqrt2$.  Since
$\bar\varepsilon-\bar\varepsilon^{-1}=24\sqrt2\ne0$, one has
\[
 p\mid V_n\quad\Longleftrightarrow\quad
 \bar\varepsilon^{\,2n}=1.
\]
Let $\rho(p)$ be the multiplicative order of
$\bar\varepsilon^2$.  Frobenius fixes $\bar\varepsilon$ when
$\chi(p)=1$ and sends it to $\bar\varepsilon^{-1}$ when
$\chi(p)=-1$.  Consequently
$\bar\varepsilon^{p-\chi(p)}=1$, and hence
$\rho(p)\mid p-\chi(p)$.

If $\rho(p)\le R$, then $p$ divides the nonzero integer
$\prod_{n\le R}V_n$.  Since $\log V_n\ll n$,
\[
 \#\{p\in\mathcal P_X:\rho(p)\le R\}\,\log X
 \le \sum_{n\le R}\log V_n\ll R^2,
\]
which is the asserted estimate.
\end{proof}

\begin{lemma}[Exceptional classes]\label{lem:qc-exceptional}
$\#\{p\in\mathcal P_X:p\mid\mathfrak U_H\}\ll H^2/\log X$.
\end{lemma}

\begin{proof}
For $H\ge\sqrt X$ this follows from the trivial estimate
$|\mathcal P_X|\ll X/\log X$.  Suppose that $H<\sqrt X$.
Then every $p\in\mathcal P_X$ exceeds~$H$ and hence divides none of
the factorials occurring in $\mathfrak U_H$.  The elementary growth
bounds $\log b_j\ll j$ and $\log V_{j+1}\ll j$ give
\[
 \log\prod_{j\le H}b_jV_{j+1}\ll H^2.
\]
Every window prime dividing $\mathfrak U_H$ therefore contributes at
least $\log X$ to the logarithm of this product, which proves the
result.
\end{proof}

\begin{hypothesis}[Transverse aligned-content mass,
$\mathrm{AC}^{\mathrm{tr}}$]\label{hyp:qc-ac}
With
\[
 \mathcal A_{\mathrm{tr}}(H)
 :=
 \sum_{\substack{a,b,c\ge2,\ a+b+c\le H\\
                   \neg(a=b=c)}}
 \log \mathcal C^*_{a,b,c},
\]
one has $\mathcal A_{\mathrm{tr}}(H)\ll H^{3+o(1)}$.
\end{hypothesis}

By~\eqref{eq:qc-aligned-divisibility}, this hypothesis is no stronger
than the earlier scalar reduced gcd-mass hypothesis
$\mathrm{GM}^{\mathrm{tr}}$: the latter implies
$\mathrm{AC}^{\mathrm{tr}}$ term by term.  The aligned formulation is the
one actually used below because it retains common-root multiplicity
without charging coincidences at unrelated roots.

\begin{remark}[Computational support]
The exact audit in \texttt{quadcorank\_verify.py} exposes the diagonal
error in the adjacent-pair formulation: at $H=20$, its diagonal mass
divided by $H^4$ is $0.187534$, whereas its total mass divided by
$H^3$ is already $3.842321$.  The na\"ive skipped-triple repair also
fails its scaling gate: its $c=a$ slice divided by $H^4$ is
$0.004833$ at $H=32$, exactly where the propagated central factors
occur.  After central deflation and removal of the one-parameter
arithmetic-progression slice, the scalar ratio is $0.005569$ at $H=20$
and $0.014732$ at $H=32$.

The independent exact audit in
\texttt{aligned\_corank\_verify.py} computes the full polynomial
$\mathcal R_{a,b,c}(T)\in\mathbb Z[T]$, not a gcd of sampled values.
It finds $\mathcal C^*_{a,b,c}=1$ for every triple through $H=14$.
At $H=20$, only two of the $675$ non-all-equal gap triples
have nontrivial reduced content, both equal to~$653$, and
$\mathcal A_{\mathrm{tr}}(20)/20^3=0.001620394$.  At $H=24$, only
$18$ of $1323$ triples are nontrivial and the corresponding ratio is
$0.007248253$.  These finite computations support, but of course do
not prove, $\mathrm{AC}^{\mathrm{tr}}$.
\end{remark}

\begin{theorem}[Quadruple corank bound]\label{thm:qc-main}
Assume $\mathrm{AC}^{\mathrm{tr}}$.  Then
$\sum_{X<p\le2X}Z(p)\ll X^{3/2+o(1)}$,
i.e., average $Z(p)\ll p^{1/2+o(1)}$.
\end{theorem}

\begin{proof}
We first record the valuation form of
Lemma~\ref{lem:qc-cert}.  Partition $[0,p-1]$ into consecutive blocks
of length at most~$H$, and suppose that one block contains the ordered
zeros $z_1<\cdots<z_m$.  Select the $m-3$ consecutive quadruples
$(z_i,z_{i+1},z_{i+2},z_{i+3})$ when $m\ge4$.  Thus the total number
$Q_p(H)$ of selected quadruples satisfies
\begin{equation}\label{eq:qc-block-lower}
 Q_p(H)\ge Z(p)-3\left\lceil\frac pH\right\rceil.
\end{equation}
Every selected quadruple has a gap triple
$\boldsymbol h=(h_1,h_2,h_3)$ with $h_i\ge2$ and
$h_1+h_2+h_3<H$.

Put $\mathcal E_H=\{p\in\mathcal P_X:p\mid\mathfrak U_H\}$.
The endpoint identity
\[
 N_h(-1)=((h-1)!)^3b_{h-1}
\]
and the definition of $\mathfrak U_H$ show that every $N_h$ with
$h\le H$ has nonzero reduction modulo every
$p\notin\mathcal E_H$.  The same holds for $N_h^\circ$: this is
immediate for odd~$h$, while for even~$h$ the factorization
$N_h=(2x+h+1)N_h^\circ$ shows that a zero reduction of $N_h^\circ$
would make $N_h$ the zero polynomial modulo~$p$.

For completeness, if two integer polynomials $f,g$, both nonzero
modulo~$p$, have $t$ distinct common roots modulo~$p$, evaluation at
those roots gives $t$ independent linear functionals annihilating the
image of the fixed, characteristic-zero-size Sylvester map.  Hence its
reduction modulo~$p$ has corank at least~$t$.  Smith normal form then
gives $v_p(\operatorname{Res}(f,g))\ge t$ whenever the integer
resultant is nonzero.  A leading-coefficient drop after reduction does
not change this fixed matrix or destroy any of the evaluation
functionals.

Fix $p\notin\mathcal E_H$, write a gap triple as $(a,b,c)$, and let
$t_p(a,b,c)$ be the number of its selected starting points.  If
$c\ne a$, every such starting point is a common root of the three
polynomials in the first line of Definition~\ref{def:qc-aligned}.
Lemma~\ref{lem:qc-aligned} gives
\[
 t_p(a,b,c)
 \le v_p(\mathcal C_{a,b,c}).
\]
If $c=a$ and $a\ne b$, let $e_p(a,b)$ be the number of those starts
that satisfy~\eqref{eq:qc-central-start}.  Lemma~\ref{lem:qc-deflation}
shows that every remaining start is a common root of the three
polynomials in the second line of Definition~\ref{def:qc-aligned}, so
the aligned-content lemma gives
\[
 t_p(a,b,a)-e_p(a,b)
 \le v_p(\mathcal C_{a,b,a}),
 \qquad
 \sum_{a,b}e_p(a,b)\le1.
\]
Finally, if $a=b=c$, the first three zeros already make every start a
common root of $N_a(x)$ and $N_a(x+a)$, whence
$t_p(a,a,a)\le v_p(S_{a,a})$.  Summing the three classes yields
\begin{equation}\label{eq:qc-valuation}
 Q_p(H)-1\le
 \sum_{\substack{a,b,c\ge2,\ a+b+c\le H\\
                  \neg(a=b=c)}}
 v_p(\mathcal C_{a,b,c})
 +\sum_{\substack{a\ge2\\3a\le H}}v_p(S_{a,a}).
\end{equation}

It remains to retain this valuation information at primes dividing the
structural carrier.  We do so by choosing the block length separately
on each zero-count level.  The primes with $Z(p)\le\sqrt X$ contribute
at most
\begin{equation}\label{eq:qc-low-level}
 \sqrt X\,|\mathcal P_X|\ll\frac{X^{3/2}}{\log X}.
\end{equation}
For dyadic $T=2^k\sqrt X$ with $T<2X$, put
\[
 \mathcal P_X(T)=\{p\in\mathcal P_X:T<Z(p)\le2T\},
 \qquad H=\left\lceil\frac{16X}{T}\right\rceil.
\]
For all sufficiently large~$X$, one has $8\le H<p$ and
\[
 3\left\lceil\frac pH\right\rceil
 \le \frac{3T}{8}+3.
\]
Equations~\eqref{eq:qc-block-lower} and~\eqref{eq:qc-valuation}
therefore imply, uniformly for
$p\in\mathcal P_X(T)\setminus\mathcal E_H$,
\begin{equation}\label{eq:qc-level-certificates}
 \frac T3\le
 \sum_{\substack{a,b,c\ge2,\ a+b+c\le H\\
                  \neg(a=b=c)}}
 v_p(\mathcal C_{a,b,c})
 +\sum_{\substack{a\ge2\\3a\le H}}v_p(S_{a,a}).
\end{equation}

For $p\notin\mathcal E_H$, saturation changes none of the
$p$-adic valuations in~\eqref{eq:qc-level-certificates}.  Moreover,
the general resultant-height estimate quoted above gives
\[
 \sum_{\substack{a\ge2\\3a\le H}}\log|S^*_{a,a}|
 \le \sum_{\substack{a\ge2\\3a\le H}}\log|S_{a,a}|
 \ll H^3\log H=H^{3+o(1)}.
\]
Summing over
$\mathcal P_X(T)\setminus\mathcal E_H$, then using
$\sum_{X<p\le2X}v_p(M)\le\log|M|/\log X$ and
$\mathrm{AC}^{\mathrm{tr}}$, gives
\[
 \frac T3\,|\mathcal P_X(T)\setminus\mathcal E_H|
 \ll\frac{H^{3+o(1)}}{\log X}.
\]
Since $Z(p)\le2T$ on this level,
\begin{equation}\label{eq:qc-generic-level}
 \sum_{p\in\mathcal P_X(T)\setminus\mathcal E_H}Z(p)
 \ll\frac{H^{3+o(1)}}{\log X}
 \ll\frac{(X/T)^{3+o(1)}}{\log X}.
\end{equation}
On the exceptional part, Lemma~\ref{lem:qc-exceptional} and the same
level upper bound give
\begin{equation}\label{eq:qc-exceptional-level}
 \sum_{p\in\mathcal P_X(T)\cap\mathcal E_H}Z(p)
 \le2T|\mathcal E_H|
 \ll\frac{TH^2}{\log X}
 \ll\frac{X^2}{T\log X}.
\end{equation}
Finally, summing~\eqref{eq:qc-generic-level} and
\eqref{eq:qc-exceptional-level} over the dyadic values
$T\ge\sqrt X$ yields convergent geometric sums:
\[
 \sum_{T\ge\sqrt X}\frac{(X/T)^{3+o(1)}}{\log X}
 +\sum_{T\ge\sqrt X}\frac{X^2}{T\log X}
 \ll X^{3/2+o(1)}.
\]
Together with~\eqref{eq:qc-low-level}, this proves the theorem.
\end{proof}

\subsubsection{A fully deflated adjacent certificate}

The aligned pencil above deliberately uses a skipped endpoint.  There is a
second, more literal certificate which keeps all three adjacent gaps and
removes the reflection-center factor from each of them.

\begin{definition}[Fully deflated adjacent content]
\label{def:qc-fd}
For $a,b,c\ge2$, put
\[
 F^{\mathrm{fd}}_{a,b,c}=N_a^\circ(x),\qquad
 G^{\mathrm{fd}}_{a,b,c}=N_b^\circ(x+a),\qquad
 J^{\mathrm{fd}}_{a,b,c}=N_c^\circ(x+a+b),
\]
and define
\[
 \mathcal R^{\mathrm{fd}}_{a,b,c}(T)
 =\operatorname{Res}_x\!\left(
 F^{\mathrm{fd}}_{a,b,c},
 G^{\mathrm{fd}}_{a,b,c}+T J^{\mathrm{fd}}_{a,b,c}
 \right),
 \qquad
 \mathcal C^{\mathrm{fd}}_{a,b,c}
 =\operatorname{cont}_T\mathcal R^{\mathrm{fd}}_{a,b,c}.
\]
Relative to height~$H$, let
$(\mathcal C^{\mathrm{fd}}_{a,b,c})^*$ denote removal of the full
prime-power part supported on~$\mathfrak U_H$.
\end{definition}

\begin{lemma}[Fully deflated quadruple certificate]
\label{lem:qc-fd}
Every polynomial $\mathcal R^{\mathrm{fd}}_{a,b,c}(T)$ is nonzero.
Let an odd prime~$p$ be fixed, partition $[0,p-1]$ into intervals,
and select sliding quadruples of consecutive zeros inside the intervals.
Apart from at most three selected quadruples in total, every selected
quadruple with gap vector $(a,b,c)$ has its start as a common root of
\[
 F^{\mathrm{fd}}_{a,b,c},\quad
 G^{\mathrm{fd}}_{a,b,c},\quad
 J^{\mathrm{fd}}_{a,b,c}\pmod p.
\]
\end{lemma}

\begin{proof}
Every root of $N_h^\circ$ is a root of~$N_h$.  The root-strip theorem
therefore places the roots of $F^{\mathrm{fd}}_{a,b,c}$ in
$-a<\operatorname{Re}x<-1$, while the roots of
$G^{\mathrm{fd}}_{a,b,c}$ lie in
$-a-b<\operatorname{Re}x<-a-1$.  These strips are disjoint, so
\[
 \mathcal R^{\mathrm{fd}}_{a,b,c}(0)
 =\operatorname{lc}(F^{\mathrm{fd}}_{a,b,c})^{(d-d_G)}
   \operatorname{Res}_x(F^{\mathrm{fd}}_{a,b,c},
                         G^{\mathrm{fd}}_{a,b,c})\ne0,
\]
where $d_G=\deg G^{\mathrm{fd}}_{a,b,c}$ and
$d=\max(d_G,\deg J^{\mathrm{fd}}_{a,b,c})$.
This proves nonvanishing of the pencil.

Let $u<v$ be one of the three adjacent zero pairs in a selected
quadruple, and put $h=v-u$.  The gap identity gives $N_h(u)=0$.
If $h$ is odd, this is already $N_h^\circ(u)=0$.  If $h$ is even,
the integral factorization
$N_h(x)=(2x+h+1)N_h^\circ(x)$ shows the same unless
\[
 2u+h+1=u+v+1\equiv0\pmod p.
\]
Because $0\le u<v\le p-1$, the exceptional congruence is equivalent
to $u+v=p-1$; call such an adjacent pair centered.  There is at most
one centered adjacent pair in any increasing zero set.  Indeed, two
different such pairs would cross: if $u_1<u_2$ and
$u_1+v_1=u_2+v_2=p-1$, then $u_1<u_2<v_2<v_1$, contradicting the
adjacency of $u_1,v_1$.  A fixed
adjacent pair occurs in at most three sliding quadruples, according as
it is their first, second, or third pair.  Thus at most three selected
quadruples are lost.  Every other selected quadruple survives all
three deflations and gives the asserted common root.
\end{proof}

\begin{hypothesis}[Fully deflated adjacent-content mass,
$\mathrm{FDAC}^{\mathrm{tr}}$]
\label{hyp:qc-fdac}
With
\[
 \mathcal A_{\mathrm{fd}}(H)
 :=\sum_{\substack{a,b,c\ge2,\ a+b+c\le H\\
                    \neg(a=b=c)}}
 \log (\mathcal C^{\mathrm{fd}}_{a,b,c})^*,
\]
one has $\mathcal A_{\mathrm{fd}}(H)\ll H^{3+o(1)}$.
\end{hypothesis}

\begin{theorem}[Fully deflated quadruple corank bound]
\label{thm:qc-fd-main}
Assume $\mathrm{FDAC}^{\mathrm{tr}}$.  Then
$\sum_{X<p\le2X}Z(p)\ll X^{3/2+o(1)}$.
\end{theorem}

\begin{proof}
Use the same block selection as in the proof of
Theorem~\ref{thm:qc-main}.  For $p\notin\mathcal E_H$, all three
fully deflated polynomials have nonzero reduction.  Lemmas
\ref{lem:qc-fd} and~\ref{lem:qc-aligned} give
\[
 Q_p(H)-3\le
 \sum_{\substack{a,b,c\ge2,\ a+b+c\le H\\
                   \neg(a=b=c)}}
 v_p(\mathcal C^{\mathrm{fd}}_{a,b,c})
 +\sum_{\substack{a\ge2\\3a\le H}}v_p(S_{a,a}).
\]
The arithmetic-progression term is unchanged.  On the dyadic level
$T<Z(p)\le2T$, the choice $H=\lceil16X/T\rceil$ and
\eqref{eq:qc-block-lower} give, for $p\le2X$,
\[
 Q_p(H)-3
 >T-3\left\lceil\frac pH\right\rceil-3
 \ge \frac{5T}{8}-6\ge\frac T2
\]
for all sufficiently large~$X$.  Thus one obtains the analogue of
\eqref{eq:qc-level-certificates}, with $T/2$ on the left.  Hypothesis
$\mathrm{FDAC}^{\mathrm{tr}}$, the unconditional
$H^{3+o(1)}$ bound for the progression resultants, and
Lemma~\ref{lem:qc-exceptional} then give exactly
\eqref{eq:qc-generic-level} and~\eqref{eq:qc-exceptional-level}.
The same dyadic summation completes the proof.
\end{proof}

\begin{remark}[Exact fully deflated audit]
\label{rem:qc-fd-computation}
The script \texttt{fully\_deflated\_corank\_verify.py} computes the
entire polynomial resultant over~$\mathbb Z[T]$.  It finds
$(\mathcal C^{\mathrm{fd}}_{a,b,c})^*=1$ for every non-progression
triple through $H=23$.  At $H=32$ there are thirteen nontrivial
records, supported at $157$, $431$, or~$653$, and
$\mathcal A_{\mathrm{fd}}(32)/32^3=0.002569880$; the exact-content
SHA-256 is
\[
 \begin{gathered}
  \texttt{437f7328a4ea7d78a62b3e781165a76a}\\
  \texttt{31f277f7ec6983f9f6c3a4e67ee05efa}.
 \end{gathered}
\]
Extending the same full-content computation to $H=36$ gives nineteen
nontrivial records, supported at $157$, $431$, $499$, $653$, or~$1297$, and
\[
 \frac{\mathcal A_{\mathrm{fd}}(36)}{36^3}=0.002595232.
\]
Its exact-content SHA-256 is
\[
 \begin{gathered}
  \texttt{196327788a7ea9adbbb141a5c6161ae9}\\
  \texttt{6125ea5b8242bc6603352080b32aac76}.
 \end{gathered}
\]
The independent script \texttt{centered\_residual\_verify.py}
reconstructs all the common roots.  The roots at $157$ and~$653$ skip
an intermediate return, and their refined return-time lists contain a
centered adjacent pair.  (One reduced content is
$24649=157^2$, not a new characteristic.)  The two roots at~$431$,
for gap vectors $(13,3,15)$ and $(15,3,13)$, are primitive and
off-center; nevertheless they are not Ap\'ery zeros, since
$\mathcal Z_{431}=\{56,374\}$.  Thus none of the height-$32$
residual roots is an actual primitive four-zero obstruction.  This
finite observation is not used in Theorem~\ref{thm:qc-fd-main}.

Primitive saturation does not remove every algebraic phantom.  The new
$p=499$ records at $H=33$ have reflected triples $(6,14,13)$ and
$(13,14,6)$, roots $417$ and~$48$, and exact return-offset lists
$(0,6,20,33)$ and $(0,13,27,33)$.  Both are primitive and off-center, but
\[
 \mathcal Z_{499}=\{67,170,203,295,328,431\},
\]
so neither root belongs to the distinguished fiber.  At $H=35$ the full
content audit also verifies
\[
 (\mathcal C^{\mathrm{fd}}_{17,14,4})^*=157,
 \qquad
 \gcd_{\mathbb F_{157}[x]}
  (F^{\mathrm{fd}}_{17,14,4},G^{\mathrm{fd}}_{17,14,4},
   J^{\mathrm{fd}}_{17,14,4})=x+119.
\]
The root $x=38$ has precisely the return offsets
$0,17,31,35$ in the selected span, so it is primitive and no adjacent
pair is centered; its reflected triple $(4,14,17)$ has root~$83$ and
offsets $(0,4,18,35)$.  Nevertheless the two independent recurrence
implementations both give $\mathcal Z_{157}=\varnothing$; the algebraic
roots are not zeros of the distinguished Ap\'ery orbit.  Thus even a
primitive common-root carrier can strictly overcount the starts needed
in the application.

The unit-scale formal short-window analogue also fails.  The exact endpoint-resultant
superset scan in \texttt{primitive\_fd\_candidate\_verify.py}, followed by
an actual three-polynomial gcd over each candidate finite field, finds at
height~$36$ the reflected pair
\[
 (a,b,c,p,x)=(5,20,10,1297,360),
 \qquad (10,20,5,1297,901).
\]
In both cases the four displayed return offsets are the complete list in
the selected span, namely $(0,5,25,35)$ and $(0,10,30,35)$ respectively,
and none of their adjacent pairs is centered.  Here
$1297\nmid\mathfrak U_{35}$ and $1297>35^2$, so these are genuinely
primitive short-span common roots outside the structural carrier.
They are still phantoms: the divided and cleared recurrences independently
give
\[
 \mathcal Z_{1297}=\{459,530,766,837\}.
\]
Among all nine saturated endpoint candidates with $p>(a+b+c)^2$ through
height~$36$, the other seven consist of three endpoint false positives and
four nonprimitive roots; none starts at an actual Ap\'ery zero.  The exact
short-range payload SHA-256 is
\[
 \begin{gathered}
  \texttt{7279c7675994a3709a2039a0b04073f8}\\
 \texttt{8e3fca346d08fe94a39f5fa4a3ee76e9}.
 \end{gathered}
\]
The endpoint superset does not miss a chain defined only by its raw return
set.  Indeed, applying the continuant congruence in
Remark~\ref{rem:qc-corank-terminal} first with $(u,v)=(a,b)$ and then with
$(u,v)=(a+b,c)$ shows that exactness of the four-return set makes the two
intervening multipliers nonzero and forces the shifted adjacent gap
polynomials to vanish.

Incremental endpoint scans in eight-span blocks extend the exact calculation
through span~$84$.  Beyond span~$36$, the blocks contain five endpoint false
positives, eight nonprimitive roots at $p=8941$, and one new reflected pair of
primitive phantoms,
\[
 (a,b,c,p,x)=(3,40,31,7411,4681),\qquad
 (31,40,3,7411,2655).
\]
Their exact return-offset lists are $(0,3,43,74)$ and $(0,31,71,74)$,
respectively; neither has a centered adjacent pair,
$7411\nmid\mathfrak U_{74}$, and the two recurrence implementations give
$\mathcal Z_{7411}=\varnothing$.  The final span block $77$--$84$ has two
tested primes above $s^2$, both endpoint false positives.  Every untested
prime already satisfies $p\le s^2$, so the combined endpoint scan is an exact
verification of \textup{(QPRS)} with $C=2$ for every gap triple of span at
most~$84$ outside the all-equal slice.  The exact block digests are hard-coded
in \texttt{primitive\_fd\_candidate\_verify.py}.

A complementary prime-first audit shows that the constant~$2$ does not
persist.  The C++ program
\texttt{primitive\_projective\_prime\_scan.cpp} computes the two independent
Ap\'ery recurrence solutions for every prime $7\le p\le500000$, groups all
indices by their point in $\mathbb P^1(\mathbb F_p)$, and tests every sliding
quadruple of consecutive occurrences in each fiber.  It checks the carrier,
center, nonwrapping, and non-all-equal conditions directly.  Among
$13{,}695{,}120$ primitive off-center windows, twenty have $p>s^2$ and four
have $p>2s^2$; all twenty belong to nondistinguished fibers.  The largest
ratio is
\[
 \frac{p}{s^2}=\frac{128047}{164^2}=4.760819\ldots,
\]
attained at
\[
 (p,x;a,b,c)=(128047,42375;41,86,37)
\]
and its reflection.  Its raw return-offset set is exactly
$(0,41,127,164)$, one has $128047\nmid\mathfrak U_{164}$, and none of the
three adjacent pairs is centered.  The independent standard-library program
\texttt{primitive\_projective\_prime\_verify.py} reconstructs all twenty
records both from projective states and from the scalar continuant recurrence.
The complete scan output has SHA-256
\[
 \begin{gathered}
  \texttt{8eeb4d6d4a6f371d8f8d87facadff137}\\
  \texttt{f66dfa476a28db1a3af8892625032bf9}.
 \end{gathered}
\]
This finite result disproves $C=2$, not the existence of some absolute
constant in \textup{(QPRS)}.

There is, however, a reason not to treat the formal support hypothesis as the
main orbit-coupling target.  In the random palindromic-orbit model, the
expected number at one prime of four-equal-state windows of span at most~$s$
has order $s^3/p^2$.  At $s=\sqrt{p/C}$, summation over $p\le P$ therefore
has heuristic order $\sqrt P/(C^{3/2}\log P)$, which diverges for every fixed
$C$.  The carrier and center exclusions have lower-order random densities.
This occupancy calculation is only a heuristic and does not disprove
\textup{(QPRS)}, but the growing formal phantoms show why a distinguished-orbit
input is structurally preferable.

Thus primitivity alone does not prove the unit-scale short-window statement:
coupling to the distinguished Ap\'ery orbit is indispensable.  The weaker
fixed-$\eta$ formulation below remains sufficient for the application.

The separate script \texttt{primitive\_chain\_census.py} audits the
generated $350{,}104$ pairs $(p,r)$ through $p\le5\cdot10^6$.  The
little-endian pair file has SHA-256
\[
 \begin{gathered}
  \texttt{5739c6e7fee4210678bc50bdda7d0a7c}\\
  \texttt{2f9fa082ab392e050556ebdc62ecac8b}.
 \end{gathered}
\]
Among $43{,}366$ sliding quadruples in the full ordered zero sets,
$1{,}418$ do not contain a centered adjacent pair, but none of those
$1{,}418$ has total span at most~$\sqrt p$.  In fact the closest observed
squared scale ratio is still the exact value
\[
 \frac{s^2}{p}=\frac{1428025}{3727}=383.156694\ldots.
\]
The unrestricted statement is false: this closest example, which is also
the first by absolute span, occurs at $p=3727$, with consecutive zeros
\[
 99<868<1011<1294
 \quad\text{and gaps}\quad (769,143,283).
\]
Thus any conjectural primitive-return strengthening must retain the
mesoscopic short-span condition used by the block argument.
\end{remark}

\begin{definition}[Short off-center window count]
\label{def:qc-short-window}
For an odd prime~$p$ and $0<\eta\le1$, write the full ordered Ap\'ery
zero set as $z_1<\cdots<z_{Z(p)}$, and let
$E_{\mathrm{sw}}(p;\eta)$ be the number of
indices $1\le i\le Z(p)-3$ such that
\[
 z_{i+3}-z_i\le\eta\sqrt p
\]
and none of the three adjacent pairs in
$z_i,z_{i+1},z_{i+2},z_{i+3}$ is centered.
\end{definition}

\begin{proposition}[Short-window reduction]
\label{prop:qc-short-window}
For every odd prime~$p$ and $0<\eta\le1$,
\[
 Z(p)<3\eta^{-1}\sqrt p+6+E_{\mathrm{sw}}(p;\eta).
\]
Consequently, for every fixed~$\eta$,
\[
 \sum_{X<p\le2X}Z(p)
 \ll_\eta \frac{X^{3/2}}{\log X}
     +\sum_{X<p\le2X}E_{\mathrm{sw}}(p;\eta).
\]
\end{proposition}

\begin{proof}
Put $q=Z(p)$ and $E=E_{\mathrm{sw}}(p;\eta)$; there is nothing to prove
when $q\le6+E$.  For
$1\le i\le q-3$, let $s_i=z_{i+3}-z_i$.  There is at most one
centered adjacent pair in the full increasing zero set, and a fixed
adjacent pair belongs to at most three sliding quadruples.  Apart from
the $E$ short off-center windows, at most three further $s_i$ can be at
most~$\eta\sqrt p$.  Hence at least $q-6-E$ of the sliding spans satisfy
$s_i>\eta\sqrt p$.  On the other hand, writing
$d_j=z_{j+1}-z_j$, each $d_j$ occurs in at most three of the sums
$s_i=d_i+d_{i+1}+d_{i+2}$, so
\[
 (q-6-E)\eta\sqrt p
 <\sum_{i=1}^{q-3}s_i
 \le3\sum_{j=1}^{q-1}d_j
 =3(z_q-z_1)<3p.
\]
This proves the pointwise inequality.  Summing it over the primes in
$(X,2X]$ and using the elementary prime-counting bound proves the
second assertion.
\end{proof}

\begin{hypothesis}[Short-window reflection, \textup{(SWR)}]
\label{hyp:qc-swr}
There is an absolute constant $0<\eta\le1$ such that
$E_{\mathrm{sw}}(p;\eta)=0$ for every odd prime~$p$; equivalently,
every four consecutive Ap\'ery zeros of total span at most
$\eta\sqrt p$ contain a centered adjacent pair.
\end{hypothesis}

\begin{corollary}[Square-root bound under short-window reflection]
\label{cor:qc-swr}
Assume \textup{(SWR)} with the fixed constant~$\eta$.  Then, for every
odd prime~$p$,
\[
 Z(p)\le
 3\left\lceil
   \frac{p}{\lfloor\eta\sqrt p\rfloor+1}
  \right\rceil+3
 <3\eta^{-1}\sqrt p+6.
\]
Consequently
$\sum_{X<p\le2X}Z(p)\ll X^{3/2}/\log X$.
\end{corollary}

\begin{proof}
The case $p=3$ is immediate from
$(b_0,b_1,b_2)\equiv(1,2,1)\pmod3$.  Suppose that $p\ge5$, put
$n=\lfloor\eta\sqrt p\rfloor$ and $H=n+1$, and partition
$\{0,\ldots,p-1\}$ into
$B=\lceil p/H\rceil$ consecutive blocks of cardinality at most~$H$.
If the blocks contain respectively $m_1,\ldots,m_B$ zeros, select all
sliding quadruples of consecutive zeros within each block.  Their total
number~$Q$ satisfies
\[
 Q=\sum_{\nu=1}^{B}(m_\nu-3)_+
 \ge Z(p)-3B.
\]
Every selected quadruple is consecutive in the full zero set and has
span at most $H-1=n\le\eta\sqrt p$, so \textup{(SWR)} makes it contain
a centered adjacent pair.  There is at most one such pair in the full zero set,
and it belongs to at most three sliding quadruples.  Hence $Q\le3$ and
\[
 Z(p)\le3B+3
 =3\left\lceil\frac{p}{n+1}\right\rceil+3.
\]
Since $n+1>\eta\sqrt p$, the ceiling is strictly less than
$\eta^{-1}\sqrt p+1$, proving the second pointwise bound.  The asserted
prime-window average follows from that bound and the standard estimate
$\#\{X<p\le2X\}\ll X/\log X$.
\end{proof}

\begin{hypothesis}[Quadratic primitive-return support,
\textup{(QPRS)}]
\label{hyp:qc-qprs}
There is an absolute constant $C\ge1$ with the following property.  Let
$p$ be an odd prime, let $a,b,c\ge2$ be not all equal, put
$s=a+b+c<p$, and let $x$ be an integer with $0\le x$ and $x+s<p$.
Define the return-offset set
\[
 \mathcal T_{p,s}(x)
 :=\{0\}\cup\{1\le h\le s:N_h(x)\equiv0\pmod p\}.
\]
If
\[
 p\nmid\mathfrak U_s,
 \qquad
 \mathcal T_{p,s}(x)=\{0,a,a+b,s\},
\]
and none of the three adjacent pairs with offsets
$0,a,a+b,s$ is centered, then
\[
 p\le C s^2.
\]
\end{hypothesis}

\begin{theorem}[Quadratic primitive support implies the square-root average]
\label{thm:qc-qprs}
Assume \textup{(QPRS)}.  Then
\[
 \sum_{X<p\le2X}Z(p)\ll_C X^{3/2}.
\]
\end{theorem}

\begin{proof}
Fix $K=64\sqrt C$.  The primes with $Z(p)\le K\sqrt X$ contribute
$O_C(X^{3/2}/\log X)$.  For dyadic $T\ge K\sqrt X$, put
\[
 \mathcal P_X(T)=\{p\in\mathcal P_X:T<Z(p)\le2T\},
 \qquad H=\left\lceil\frac{16X}{T}\right\rceil,
\]
and use the block selection from the proof of
Theorem~\ref{thm:qc-main}.  For all sufficiently large~$X$,
\[
 H\le\frac{16\sqrt X}{K}+1
 \le\frac{\sqrt X}{2\sqrt C},
 \qquad CH^2<X.
\]
Put $\mathcal E_H=\{p\in\mathcal P_X:p\mid\mathfrak U_H\}$.

Suppose that $p\in\mathcal P_X(T)$ does not divide
$\mathfrak U_H$, and consider a selected consecutive quadruple with
start~$x$ and gap vector $(a,b,c)$.  Its span $s=a+b+c$ is less
than~$H$, so $p\nmid\mathfrak U_s$.  By the projective-orbit gap
identity in Remark~\ref{rem:orbit}, because $b_x\equiv0\pmod p$ one has
\[
 N_h(x)\equiv0\pmod p
 \quad\Longleftrightarrow\quad
 b_{x+h}\equiv0\pmod p
 \qquad(0\le h\le s).
\]
The global consecutiveness of the quadruple therefore gives
$\mathcal T_{p,s}(x)=\{0,a,a+b,s\}$.  If it contained no centered
adjacent pair and $a,b,c$ were not all equal,
\textup{(QPRS)} would imply
\[
 p\le Cs^2<CH^2<X<p,
\]
a contradiction.

Let $A_p(H)$ be the number of selected quadruples with $a=b=c$.
The preceding paragraph and the uniqueness of the centered adjacent pair
give
\[
 Q_p(H)\le A_p(H)+3.
\]
Every progression start counted by $A_p(H)$ is a common root of
$N_a(x)$ and $N_a(x+a)$ for some $3a<H$.  The fixed Sylvester-corank
argument therefore gives
\[
 A_p(H)\le\sum_{\substack{a\ge2\\3a<H}}v_p(S_{a,a}).
\]

On the other hand,
\[
 Q_p(H)\ge Z(p)-3\left\lceil\frac pH\right\rceil
 >T-3\left(\frac T8+1\right)
 =\frac{5T}{8}-3.
\]
Consequently, for all sufficiently large~$X$ and every
$p\in\mathcal P_X(T)\setminus\mathcal E_H$,
\[
 \frac T2\le\sum_{\substack{a\ge2\\3a<H}}v_p(S_{a,a}).
\]
Summing over these generic primes and using
$\sum_{3a<H}\log|S_{a,a}|\ll H^3\log H$ gives
\[
 \sum_{p\in\mathcal P_X(T)\setminus\mathcal E_H}Z(p)
 \le2T\,|\mathcal P_X(T)\setminus\mathcal E_H|
 \ll\frac{H^3\log H}{\log X}
 \ll H^3.
\]
On the exceptional part, Lemma~\ref{lem:qc-exceptional} gives
\[
 \sum_{p\in\mathcal P_X(T)\cap\mathcal E_H}Z(p)
 \le2T\,\#\{p\in\mathcal P_X:p\mid\mathfrak U_H\}
 \ll\frac{TH^2}{\log X}
 \ll\frac{X^2}{T\log X}.
\]
Since $H\ll X/T$, summing both geometric series over
$T\ge K\sqrt X$ and adding the low range proves the theorem.
\end{proof}

\begin{corollary}[Moment threshold $k=5$]
Under any of $\mathrm{AC}^{\mathrm{tr}}$,
$\mathrm{FDAC}^{\mathrm{tr}}$, \textup{(SWR)}, or \textup{(QPRS)},
$(\mathrm{HM})_5$ implies
$G_n=e^{o(n)}$ for every~$n$.
\end{corollary}

\begin{proof}
Theorem~\ref{thm:qc-main}, Theorem~\ref{thm:qc-fd-main},
Theorem~\ref{thm:qc-qprs}, or Corollary~\ref{cor:qc-swr}, as appropriate,
gives
$\lambda_X\le X^{-1}\sum_{p\in\mathcal P_X}Z(p)
\ll X^{1/2+o(1)}$.  Repeating the proof of
Theorem~\ref{thm:hm-pointwise} with this improved estimate and $k=5$
gives
\[
 \max_{m<X^2}K_X(m)\ll X^{2/5+o(1)}\lambda_X+1
 \ll X^{9/10+o(1)}.
\]
The same dyadic decomposition used there now yields
$\log G_n\ll n^{9/10+o(1)}=o(n)$.
\end{proof}

\begin{remark}[The corank ladder terminates at $1/2$]
\label{rem:qc-corank-terminal}
Within a single triple of zeros the certificates are
algebraically dependent:
$N_{h_1+h_2}(x)\equiv N_{h_2}(x{+}h_1)\,N_{h_1+1}(x)
\pmod{N_{h_1}(x)}$,
so the triple level cannot improve beyond $3/5$.
The quadruple is the first configuration supplying two potentially
transverse certificates.  The adjacent self-gcd, the propagated
palindromic center, and the one-parameter arithmetic-progression slice
must all be removed as above.  The remaining input can be formulated
either as the aligned-content estimate $\mathrm{AC}^{\mathrm{tr}}$ or
as the fully deflated adjacent estimate
$\mathrm{FDAC}^{\mathrm{tr}}$; the older scalar gcd-mass estimate is a
sufficient but unnecessarily strong version of the first formulation.
The short-window condition \textup{(SWR)} and the quadratic support condition
\textup{(QPRS)} instead bypass aggregate content mass by ruling out generic
primitive windows at a suitably chosen block scale.
Higher tuples give strictly weaker exponents within the same corank
bookkeeping.
\end{remark}

\subsection{The parity barrier and the central open problem}

\begin{openproblem}[Parity barrier]\label{op:parity}
Find a map $r\mapsto\mathfrak G_r$ assigning to every integer $r\ge2$ a
\emph{nonzero} integer $\mathfrak G_r$, defined by $r$ alone, such that
\begin{enumerate}
\item[\textup{(D)}] \textbf{Detection:} for every prime $p>r$,
  $p\mid b_r\;\Longrightarrow\;p\mid\mathfrak G_r$;
\item[\textup{(H)}] \textbf{Height:} $\log|\mathfrak G_r|\le r^{o(1)}$.
\end{enumerate}
The sequence $\mathfrak G_r=b_r$ satisfies \textup{(D)} with
$\log b_r\sim3.5255\,r$: the entire content is the height condition.
\end{openproblem}

\begin{remark}[What a solution yields]
With $(\log r)^{O(1)}$ heights,
summing $\omega_{(X,2X]}(\mathfrak G_r)$
over $r<2X$ gives average $Z(p)\ll p^{o(1)}$;
the exponent~$c$ drops to~$o(1)$ and the full
conjecture follows from $(\mathrm{HM})_3$ alone.
A solution would \emph{not} by itself prove $(\mathrm{AT}'')$:
the atom count is a diagonal selection,
and certificates control the total divisor mass,
not its concentration on the diagonal.
It would, however, breach the parity barrier:
every certificate in this paper requires
at least two vanishing conditions at a single prime.
\end{remark}

\begin{remark}[Known failure modes]
\textup{(i)} The Fermat detector $1-b_r^{p-1}$ has degree growing
with~$p$.
\textup{(ii)} The Cartier/Hasse--Witt operator is $p$-semilinear;
its entries lift to no fixed integer.
\textup{(iii)} Character-sum realizations detect residues
modulo~$p{-}1$, not modulo~$p$.
\textup{(iv)} Weight: twisted traces have weight~$\ge2$, so
$p\mid S$ does not force $S=0$.
\textup{(v)} The digit ladder (variation of parameters at
depth~$\ge4$) is parity-neutral: each deeper congruence introduces
one relation and one new local unknown, with zero net global
information gain.
\textup{(vi)} Tensor constructions (symmetric powers) dilute parity
rather than breaking it.
\end{remark}

\begin{proposition}[Diagonal transport identity]
\label{prop:diagonal-transport}
For all $h \ge 1$,
\[
  N_{h+1}(x{-}1) = P(x)\,N_h(x) - (x{+}1)^6\,N_{h-1}(x{+}1).
\]
Equivalently, the gap-polynomial family satisfies a
\emph{shift relation}: the polynomial triples
$(N_{h+1}(x{-}1),\, N_h(x),\, N_{h-1}(x{+}1))$ are connected by the
\emph{same} coefficients $P(x)$ and $(x{+}1)^6$ independently of~$h$.
\end{proposition}

\begin{proof}
The standard $h$-recurrence expands
$\det D_{h+1}(x)$ along the last row;
the present identity expands $\det D_{h+1}(x{-}1)$
along the \emph{first} row.
\end{proof}
